\documentclass{article}\begin{document}For an indexed family of actual spans, chosen component lifts and an ordinary payload, forwardUnit and forwardComp prove the identity and composite base diagrams for \#198's assembled forward comparison between the tagged weighted pulls and the chosen weighted pull of the single coproduct span. Their equations use the inverse component of the actual SpanPull.forward isomorphism, in the orientation from pulling then comparing to comparing then pulling. UnitTag first identifies the chosen identity-route comparison under tagging using the factorization through the identity lift. CompTag cancels against the actual composite cartesian route and uses both functoriality of tagging and factorization of the component comparison. No route unit or composition correction is identified with an identity. The whole and part constraints are the whiskered Route.unitIso and Route.compIso, and the tagged constraints apply these maps at their own tags. The proof transfers the component diagrams through the strict inverse fibre presentations, and then through compareTag using actual functor equalities at each base object. ForwardApp computes the comparison at a tagged object with its endpoint correction. This is coherence of the forward comparison of weighted coproduct fibres for all tags; it does not yet establish those diagrams for the final All.evaluationBase, which uses the inverse mate and a second (tagged evaluation) base square. Further obligations are to transfer the proved diagrams across the chosen inverse equivalences retaining their units and counits, check the zero and positive tagged evaluation constraints, and paste. Mixed compatibility with payload change, bundled higher coherence and flattening remain open. The generic result does not add any split laws or change the existing concrete all-length universe restriction. Consult pinned manuscript 458ff36a, lines 929--953; the distinction between the two forward and inverse comparisons is in SpanPull.forward and SpanPull.All.evaluationBase.\end{document}
